% Generated draft: requires mathematical and source review.
\documentclass[11pt]{article}
\usepackage[margin=25mm]{geometry}
\usepackage{fontspec}
\setmainfont{Latin Modern Roman}
\usepackage{amsmath,amssymb,mathtools,unicode-math}
\setmathfont{Latin Modern Math}
\usepackage{enumitem,needspace,xurl,hyperref}
\setlength{\parindent}{0pt}
\setlength{\parskip}{5pt}
\setlength{\emergencystretch}{3em}
\newcommand{\N}{\mathbb{N}}
\newcommand{\Z}{\mathbb{Z}}
\newcommand{\Q}{\mathbb{Q}}
\newcommand{\R}{\mathbb{R}}
\newcommand{\C}{\mathbb{C}}
\newcommand{\F}{\mathbb{F}}
\newcommand{\A}{\mathbb{A}}
\newcommand{\PP}{\mathbb{P}}
\newcommand{\E}{\mathbb{E}}
\newcommand{\eps}{\varepsilon}
\newcommand{\dd}{\,\mathrm{d}}
\begin{document}
\section*{UnsolvedMath ID: 30001408}
\textbf{Source problem code:} catalogue entry
\subsection{Definitions and mathematical statement}
For a positive integer \(n\), let
\[
K_0^n=\{K\subset\R^n:K\text{ is compact and convex, and }0\in\operatorname{int}_{\R^n}K\}.
\]
Write \(E=\mathbb R_+\cup\{\infty\}\), retaining the source's notation \(\mathbb R_+\): the supplied record does not specify whether zero belongs to this set.

A valuation with values in \(E\) is a map \(\Phi:K_0^n\to E\) satisfying
\[
\Phi(K\cup L)+\Phi(K\cap L)=\Phi(K)+\Phi(L)
\]
for all \(K,L\in K_0^n\) for which \(K\cup L\) is convex. In this case both the union and intersection belong to \(K_0^n\). Addition is interpreted in the extended nonnegative range, with \(a+\infty=\infty\), including \(a=\infty\); the identity does not permit cancellation of infinite terms.

Let
\[
\mathrm{SL}(n,\R)=\{A\in\R^{n\times n}:\det A=1\},\qquad AK=\{Ax:x\in K\}.
\]
Invariance means
\[
\Phi(AK)=\Phi(K)\qquad(A\in\mathrm{SL}(n,\R),\ K\in K_0^n).
\]
Under the usual interpretation of homogeneity of real degree \(q\), the condition is
\[
\Phi(tK)=t^q\Phi(K)\qquad(t>0,\ K\in K_0^n),
\]
where \(tK=\{tx:x\in K\}\) and \(t^q\infty=\infty\). The source specifies neither a degree nor an allowed set of degrees, and does not clarify whether the classification is to be stated separately for each degree or across all degrees.

Under the usual Hausdorff-topology interpretation, upper semicontinuity means that for every \(K\in K_0^n\) and every sequence \((K_j)_{j\geq1}\subset K_0^n\),
\[
d_H(K_j,K)\longrightarrow0
\quad\Longrightarrow\quad
\limsup_{j\to\infty}\Phi(K_j)\leq\Phi(K),
\]
where \(d_H\) is the Hausdorff distance induced by the Euclidean norm and the inequality uses the extended order. The supplied record does not explicitly name this topology.

\textbf{Classification problem.} Characterize all upper-semicontinuous, \(\mathrm{SL}(n,\R)\)-invariant, homogeneous valuations
\[
\Phi:K_0^n\longrightarrow\mathbb R_+\cup\{\infty\}.
\]
The value \(\infty\) is expressly allowed. No finiteness or properness condition is supplied, and none is imposed here. The source does not state the admissible dimensions; the definition above describes the domain for positive integer \(n\), without resolving whether the intended problem excludes any such dimensions. These dimension, degree, range, and topology qualifications require confirmation before a unique fully specified classification target can be fixed.
\subsection{Short English statement}
Classify homogeneous valuations on compact convex bodies whose interiors contain the origin, assuming upper semicontinuity and invariance under determinant-one linear maps. Infinite values are explicitly permitted. The intended dimensions, homogeneity degrees, convention for \(\mathbb R_+\), and topology require confirmation from the source.
\subsection{Sources}
\begin{itemize}
\item Supplied local UnsolvedMath registry record, permanent ID \(30001408\), title: \emph{Invariant Homogeneous Valuations on Convex Bodies}; exact source problem code: \texttt{catalogue entry}. Collection: Oberwolfach Reports Open Problems. The supplied proposed year is 2010; no individual proposer is named.
\item \emph{Mini-Workshop: Valuations and Integral Geometry} (2010), Oberwolfach Reports, pp.\ 141--178, DOI \texttt{10.4171/owr/2010/04}. Supplied source and literature reference, labelled ``original Oberwolfach problem report'': \url{https://doi.org/10.4171/owr/2010/04}.
\item Supplied literature reference, labelled ``Ludwig and Reitzner'': \url{https://annals.math.princeton.edu/wp-content/uploads/annals-v172-n2-p09-p.pdf}. No title or further bibliographic details are supplied in the record.
\item Literature assessment dated \texttt{2026-08-21}: the supplied curation labels the problem open and states that the report leaves the extended-range classification open, while a later real-valued classification does not settle that extension. This is supplied dated curation, not an independently verified assessment; time-sensitive status and cited claims require independent verification.
\item The supplied statement status is \texttt{corrected\_verified}; its accompanying verification note describes enclosing TeX notation in mathematical delimiters without changing the mathematical content. This metadata does not establish mathematical correctness or resolve the formulation qualifications above.
\end{itemize}
\end{document}
